\chapter{Introduction}
\label{chap:Intro}

Technology, like most things, evolves, and with it, so do all related disciplines, 
such as science and engineering, among others. And, in turn, the teaching of these 
disciplines must also evolve.

This thesis presents a QEMU-based emulation of the STM32F405 microcontroller 
and the LIS3DH accelerometer, providing a virtual platform that runs unmodified 
firmware and replicates the behavior of real devices. The intention with this virtual 
environment is to replace some of the physical laboratory tools used in the 
Master's Program in Electronic Systems Engineering at the Polytechnic University of 
Madrid, allowing students to experiment with microcontroller peripherals and I²C-based 
sensors using only their personal computers, while preserving the pedagogical 
value of working closely with the hardware.


%Emulation is not a new concept. In computer engineering, emulation means that one 
%system (the host) behaves like another (the guest), so that the latter's software 
%and peripherals function without modification. This is the standard definition of 
%emulation in computing \cite{1}. However, when applied specifically to hardware, the 
%goal is to reproduce registers, busses, timing behavior, and I/O interfaces so that 
%the emulated device can be used in place of the original hardware system.
%
%Hardware emulation emerged in the second half of the 1980s, as verification tools. 
%This hardware emulation, better known as hardware emulation on-chip, works by 
%programming FPGA-based hardware using a hardware description language 
%(such as Verilog or VHDL) of a chip or SoC.
%
%On the other hand, software-based hardware emulation is performed purely in software, 
%without requiring any additional hardware. In this type of emulation, a 
%program creates a virtual hardware environment that mimics processors, memory 
%maps, and devices, so that firmware or operating systems compiled for the target 
%board run without modification on the host.
%
%This thesis presents a QEMU-based emulation of the STM32F405 microcontroller 
%and the LIS3DH accelerometer, providing a virtual platform that executes 
%unmodified firmware and reproduces the behavior of real devices. 
%The proposed environment aims to replace part of the physical laboratory setup, enabling students 
%to experiment with microcontroller peripherals and I2C-based sensing using only 
%their personal computers, while preserving the pedagogical value of working close to the hardware.

%==== Motivation ====%
\section{Motivation}
\label{chap:Motivation}

%==== Objectives ====%
\section{Objectives} 
\label{chap:Objectives}

%==== Overview ====%
\section{Overview}
\label{chap:Overview}